\section{Objetivo de RLHF, reward model y DPO: fórmulas parecidas, procesos de optimización distintos}

Con la retroalimentación de la evaluación disponible, esta sección pasa al entrenamiento con preferencias. RLHF primero entrena un reward model con comparaciones humanas y después optimiza la policy; DPO convierte el mismo tipo de señal de preferencia en una pérdida directa de tipo clasificación. Ambos buscan aumentar la probabilidad de la chosen response frente a la rejected response, pero difieren en el muestreo de datos, el rollout en línea, la estabilidad y el costo de cómputo.

\subsection{RLHF objective: recompensa y restricción de KL}

Esta subsección empieza por el objetivo estándar. Si solo se maximiza la learned reward, la policy puede encontrar fallas del reward model y alejarse del base behavior; por eso se añade la KL divergence respecto al reference model, que limita la deriva de la distribución. PPO es el optimizador habitual, pero el objetivo en sí no equivale a PPO.

\lecturefigure{slide-059.jpg}{Notación del RLHF objective: policy, reference, prompt, response y reward}{official deck, p.~59}

\lecturefigure{slide-060.jpg}{Primera parte del RLHF objective: maximizar la puntuación del reward model}{official deck, p.~60}

\lecturefigure{slide-061.jpg}{Segunda parte del RLHF objective: usar KL para impedir que la policy se aleje de la reference}{official deck, p.~61}

El objetivo se escribe como
\[
\max_{\theta}\;\mathbb{E}_{x\sim D,\,y\sim\pi_{\theta}(\cdot\mid x)}
\left[r_{\phi}(x,y)-\beta\log\frac{\pi_{\theta}(y\mid x)}{\pi_{\mathrm{ref}}(y\mid x)}\right].
\]
donde $x$ es el prompt, $y$ es la respuesta muestreada por la policy $\pi_{\theta}$, $r_{\phi}$ es el reward model, $\pi_{\mathrm{ref}}$ es el reference model y $\beta$ controla la KL penalty. Si $\beta$ es demasiado pequeño, aparecen la deriva y el reward hacking; si es demasiado grande, resulta difícil cambiar el comportamiento.

\teachervoice{00:40:00--00:41:55, el ponente descompone la fórmula en dos fuerzas, «perseguir la reward» y «no alejarse demasiado de la base», y enfatiza que la reference constraint no es un adorno, sino lo que impide que la learned reward se explote en exceso.}

\subsection{Preference modeling: por qué se usa la pairwise comparison}

A una persona le cuesta asignar de forma estable una puntuación absoluta de 0--10 a una respuesta abierta, pero le resulta más fácil juzgar cuál es mejor entre A y B. Por eso el reward model suele usar chosen/rejected pairs y aprende un orden relativo. La calidad de los datos depende del prompt, de la forma de generar los candidatos, de la guía de anotación y del annotator agreement.

\lecturefigure{slide-062.jpg}{Preference reward modeling: la comparación relativa supera a la supervisión directa de puntuaciones absolutas}{official deck, p.~62}

El modelo de Bradley--Terry se escribe como
\[
P(y_w\succ y_l\mid x)=\sigma\!\left(r_{\phi}(x,y_w)-r_{\phi}(x,y_l)\right),
\]
donde $y_w$ es la preferred/chosen response, $y_l$ es la rejected response, $r_{\phi}$ es una reward escalar y $\sigma$ es la sigmoide. El entrenamiento minimiza la log-probabilidad negativa de que la chosen gane.

\begin{warningbox}{Los preference data no equivalen a una verdad objetiva}
Las anotaciones reflejan una población concreta, unas políticas, una interfaz y un conjunto de candidatos. Si ambos candidatos son malos, la pairwise label de todos modos elige uno; si la diferencia entre candidatos es demasiado evidente, el reward model no aprende las fronteras finas. Conviene conservar tie/skip, la información de los anotadores y el disagreement, en lugar de guardar solo etiquetas binarias.
\end{warningbox}

\teachervoice{00:42:00--00:43:44, Nathan explica que supervisar directamente una scalar reward no da buenos resultados, mientras que la pairwise preference es más natural. El punto del curso no es que las comparaciones humanas carezcan de sesgo, sino que el juicio relativo suele ser más estable que la puntuación absoluta.}

\subsection{De la «gradient ascent directa» a DPO}

Con un pairwise reward objective, cabe preguntarse si es indispensable entrenar explícitamente un reward model y después ejecutar RL. DPO aprovecha el óptimo KL-regularized, escribe la reward implícita como el log-ratio entre policy y reference, y aplica directamente una logistic classification a cada chosen/rejected pair.

\lecturefigure{slide-063.jpg}{Pregunta: ¿se puede aplicar gradient ascent directamente al preference objective?}{official deck, p.~63}

\lecturefigure{slide-064.jpg}{Respuesta: Direct Preference Optimization reescribe el objetivo como una pérdida de preferencia directa}{official deck, p.~64}

La DPO loss es
\[
\mathcal{L}_{\mathrm{DPO}}(\theta)=-\mathbb{E}\log\sigma\!\left(
\beta\left[
\log\frac{\pi_{\theta}(y_w\mid x)}{\pi_{\mathrm{ref}}(y_w\mid x)}-
\log\frac{\pi_{\theta}(y_l\mid x)}{\pi_{\mathrm{ref}}(y_l\mid x)}
\right]\right).
\]
donde $\pi_{\theta}$ es la policy que se entrena, $\pi_{\mathrm{ref}}$ es la reference, $y_w,y_l$ son la chosen y la rejected, y $\beta$ controla la intensidad con que se permite alejarse de la reference. DPO aumenta el log-ratio relativo de la chosen y, al mismo tiempo, reduce el log-ratio relativo de la rejected.

\teachervoice{00:43:44--00:46:20, el ponente resume el atractivo de DPO en una implementación mínima que no requiere un RL rollout explícito; pero enseguida enfatiza que DPO y PPO son optimizers muy distintos: que la derivación de sus objetivos esté conectada no significa que su dinámica de entrenamiento sea la misma.}

\subsection{DPO core facts y diferencias con PPO}

Esta subsección separa la comodidad de implementación de los límites de capacidad. DPO solo necesita preference pairs fuera de línea y su entrenamiento se parece a una supervised classification; PPO muestrea en línea a partir de la policy actual, usa la reward para estimar la advantage y realiza actualizaciones de varios pasos mediante clipping/regularization. PPO es más complejo y más costoso, y también puede extraer más de la reward signal.

\lecturefigure{slide-065.jpg}{DPO core facts: implementación sencilla, adopción rápida, pero con supuestos}{official deck, p.~65}

\lecturefigure{slide-066.jpg}{DPO y PPO/REINFORCE son optimizadores y ciclos de datos distintos}{official deck, p.~66}

\begin{knowledgebox}{Términos clave: objective y optimizer}
\begin{tabularx}{\textwidth}{L{0.18\textwidth} X}
\toprule
Concepto & Significado \\
\midrule
Preference objective & Busca que la chosen sea más probable que la rejected y, a la vez, limita cuánto se aleja de la reference. \\
DPO & Optimiza directamente el log-ratio policy/reference con pairs fijos fuera de línea. \\
PPO & Hace rollout desde la policy actual y realiza actualizaciones on-policy restringidas con base en la reward/advantage. \\
REINFORCE & Policy-gradient estimator clásico que multiplica el retorno muestreado por el score-function gradient. \\
\bottomrule
\end{tabularx}
Un mismo objetivo de alto nivel puede aproximarse con distintos optimizers; la estabilidad, la eficiencia de muestras y el comportamiento final no tienen por qué coincidir.
\end{knowledgebox}

\teachervoice{En la Q\&A, 01:09:02--01:09:40, Nathan ofrece un juicio basado en la experiencia: si PPO se logra ajustar bien, a menudo extrae un poco más de rendimiento de los mismos datos; pero la ventaja se ubica en fine margins y su costo es una gran cantidad de modelos e iteraciones de entrenamiento.}

\subsection{Resumen de la sección}

El núcleo de RLHF es la learned reward junto con la reference constraint; la pairwise preference es la forma de datos habitual del reward modeling; DPO reescribe la reward implícita como una pérdida directa de clasificación de preferencias. La conexión matemática no borra las diferencias entre optimizers. Para evaluar DPO/PPO hay que considerar el release final, los datos, el cómputo, la estabilidad y las preferencias humanas, no solo la extensión de la implementación.
